\chapter{Background and Related Work}
\label{chap:related-work}

This chapter reviews the work this thesis builds on and positions its
contribution against it. It covers four areas: measuring notebook
reproducibility, restoring it automatically, using large language models
(LLMs) for dependency resolution and program repair, and the FAIR data
principles that motivate recording repair outcomes in a knowledge graph. The
chapter closes by stating, precisely, the gap none of this prior work closes
on its own.

\section{Reproducibility of Computational Notebooks}
\label{sec:rw-reproducibility}

Computational reproducibility means that a published analysis can be
re-executed by someone else and produce the same result. For a Jupyter
notebook, the weaker and more basic question is whether it can be
\emph{re-executed at all}: does the code still run, independent of whether
the numbers come out identical. This thesis targets that weaker property,
re-executability, as the necessary first step before reproducibility in the
fuller sense can even be discussed.

A large-scale study by Pimentel et al.~\cite{pimentel2019msr} measured
this directly. Out of 1.4~million notebooks collected from GitHub, only 24.11\% of
the notebooks they attempted ran without an error, and only 4.03\% produced
results matching those already stored in the notebook. Samuel and
Mietchen~\cite{samuel2024gigascience} repeated this measurement for
notebooks linked from biomedical publications on PubMed Central and found
the same low reproducibility rate in that domain. This is the dataset this
thesis' own working data is derived from, described in
\Cref{chap:problem-dataset}. Both studies agree on the dominant cause: most
failures are not bugs in the analysis logic, but missing or incompatible
software dependencies.

Detecting these failures is itself useful. Samuel and
K{\"o}nig-Ries~\cite{samuel2020reproducemegit} built ReproduceMeGit, a
visualisation tool that runs the notebooks in a GitHub repository and shows,
for each one, whether it succeeded, crashed, or produced different output
than before. Tools like this make the scale of the problem visible to a
repository owner, but they stop at reporting: they do not explain a failure
in terms a non-programmer can act on, and they do not attempt to fix
anything.

\section{Restoring Executability}
\label{sec:rw-restoring}

A separate line of work tries to go one step further and automatically
restore a broken environment rather than only reporting that it is broken.

\subsection{Static Analysis and Knowledge Graphs}

Wang et al.~\cite{wang2020ase} studied how execution environments could
be reconstructed by analysing a notebook's own code. They followed this
with SnifferDog~\cite{wang2021snifferdog}, a tool that statically
inspects a notebook's imports and matches them against a database of
known package APIs. SnifferDog correctly inferred the required
environment for 92.6\% of the notebooks it was tested on. This result
makes an important point that motivates the present thesis: a notebook
that looks broken is often not broken in its logic at all. It simply
lacks the right dependency information, which can in many cases be
recovered automatically.

A separate family of tools builds this dependency knowledge as an explicit
knowledge graph rather than analysing one notebook at a time. Horton and
Parnin's DockerizeMe~\cite{horton2019dockerizeme} introduced this idea: it encodes
package-to-module mappings from Libraries.io in a graph database and uses it
to infer, and then containerise, the dependencies a Python snippet needs.
Ye et al.'s PyEGo~\cite{ye2022pyego} extended the same idea to a much larger
graph covering third-party packages, the Python interpreter version, and
system libraries together, reporting a 0.4--3.5$\times$ accuracy improvement
over DockerizeMe. Cheng et al. added explicit conflict detection between
transitive dependencies with PyCRE~\cite{cheng2022pycre}, and later refined
the approach with ReadPyE~\cite{cheng2024readpye}, which reached 79.75\%
successful environment reconstruction, the strongest result among the
knowledge-graph approaches reviewed here.

These knowledge-graph approaches share one structural limitation: they can
only recommend a package or version that is represented in their indexed
data. Such a graph must be maintained and updated as new PyPI releases
appear, and a package or release that has not yet been indexed is
unavailable to it until the graph is rebuilt. This limitation motivates
replacing a fixed package-and-version snapshot with a live query against
the package registry itself, which is the design this thesis adopts
(\Cref{chap:architecture}).

\subsection{LLMs for Dependency Resolution}
\label{sec:rw-llm-dep}

Bartlett et al.'s PLLM~\cite{bartlett2025pllm} is the closest prior
system to this thesis and its main point of comparison. PLLM uses an
LLM, retrieval-augmented with live PyPI data, to resolve missing or
incompatible Python dependencies from execution feedback. In their
experiments, an LLM using retrieval-augmented generation (RAG) against
live package data outperformed the strongest knowledge-graph baselines
by a wide margin (+15.97\% fixes over ReadPyE, +21.58\% over PyEGo).
Results were especially strong on notebooks that depend on
machine-learning libraries -- relevant here, since the dataset in
\Cref{chap:problem-dataset} is dominated by exactly such packages
(\texttt{scikit-learn}, \texttt{scipy}, \texttt{tensorflow}). Gao et
al.'s RAG survey~\cite{gao2024ragsurvey} and Tao et al.'s survey of
retrieval-augmented code generation~\cite{tao2025racg} both confirm
that grounding a generative model's output in retrieved, verifiable
facts, rather than the model's own memorised knowledge, is established
practice for exactly this kind of task. This is the basis for Lewis et
al.'s original formulation of retrieval-augmented
generation~\cite{lewis2020rag}.

PLLM differs from this thesis in four ways that matter for the present
design: it operates on standalone Python files, not on notebooks embedded in
a reproducibility pipeline; it uses unbounded multi-round iteration rather
than a fixed number of repair attempts; it produces no human-readable
explanation of the failure; and it does not feed its results into any
knowledge graph. \Cref{sec:rw-gap} returns to each of these points directly.

\subsection{LLM-Based Program Repair and Notebook-Specific Work}
\label{sec:rw-apr}

More broadly, LLM-based automated program repair (APR) is an active
field. Yang et al.'s survey~\cite{yang2025aprsurvey} organises 63
systems published between 2022 and 2025 into four design paradigms:
fine-tuning, prompting, procedural pipelines, and agentic frameworks.
Agentic frameworks, where the model plans and re-plans its own next
action, perform well on complex, multi-file bugs, but at a real cost in
latency and control. That trade-off directly motivates a design
decision explained in \Cref{chap:architecture}: this thesis structures
its repair layer as a fixed, procedural pipeline, not an autonomous
agent, because the target workload is hundreds of independent, mostly
single-dependency notebook failures processed in batch, not a small
number of complex interactive sessions. The same survey also confirms
that retrieval-augmented context improves repair quality across every
paradigm it studied, reinforcing the PyPI-grounding decision from
\Cref{sec:rw-llm-dep}.

Two lines of work are specific to notebooks. Grotov et
al.~\cite{grotov2024untangling} released a dataset of buggy notebooks,
together with a proposed agent design. They followed it
with a full implementation~\cite{grotov2024debug} that observes the
output of each execution attempt, decides what to change, applies it,
and tries again. That second paper establishes a point this thesis
relies on directly: because a notebook keeps state between cells, a
candidate fix can only be judged correct by re-running the
\emph{whole} notebook from the top, not just the cell that originally
failed (\Cref{chap:architecture} adopts this rule for the same reason).
Xia and Zhang's ChatRepair~\cite{xia2024chatrepair} similarly feeds
re-execution feedback back into the next repair attempt, in a
conversational loop, achieving 162 of 337 fixes on a general Python bug
benchmark. The same feedback principle underlies the bounded two-round
repair mode built for this thesis (\Cref{sec:orchestrator}), with one
difference: the budget here is fixed at two rounds rather than left open.
Reyes et
al.'s Byam~\cite{reyes2025byam} and Tang et al.'s
SynFix~\cite{tang2025synfix} both target dependency-aware repair
outside the notebook setting: breaking dependency updates in Java, and
structured cross-file repair, respectively. Both motivate giving the
repair model structured dependency context rather than the bare error
message. This thesis follows that principle in a different form. The
repair prompt carries the classified error (type, message, failing
module, subtype, root-cause hint), the failing cell's source where it
could be recovered, and release evidence retrieved live from PyPI; a
repository's declared requirements are used later, when FixApplicator
reconstructs the execution environment, not as evidence handed to the
repair model (\Cref{sec:m-repair-nodeps}).

On the interactive side, Jupyter AI~\cite{jupyterai2023} lets a user ask an
LLM inside JupyterLab to explain an error or generate a fix. It is a useful
point of contrast: Jupyter AI is designed for one user working on one
notebook at a time and must be triggered manually after every failure. It
has no route into a batch reproducibility pipeline and no connection to the
structured metadata such a pipeline already records. The system built for
this thesis is designed for exactly the opposite situation: no user present,
hundreds of notebooks, and existing structured failure metadata to build on.

Model and prompt choice also draw on this literature. Szalontai et
al.~\cite{szalontai2024codellama} find that a smaller, code-specialised
model can outperform a larger general-purpose one on bug-fixing tasks,
and that prompt design affects repair quality about as much as model
size does. A broader empirical study of LLM-based
APR~\cite{empiricalapr2025} confirms that code-specialised models tend
to lead on Python repair specifically, across more than 600,000
patches. Dinh et al.'s comparison of prompting techniques for code
generation~\cite{dinh2024codeprompteval} informs the structured-output,
few-shot prompt design in \Cref{chap:methodology}. Chain-of-thought
prompting~\cite{wei2022cot} was considered and not adopted: the prompts
request a structured JSON answer with few-shot examples rather than
explicit intermediate reasoning, because the benefit of additional
reasoning prompts for locally deployable models of this size is
uncertain and a second prompt technique would have made the comparison
less controlled. Dinh et al.
in particular find that including a project's declared package list in
the prompt measurably improves dependency-aware generation. The repair
prompt in this thesis applies the same idea to a different, verifiable
source: instead of the project's own package list, it supplies the
candidate releases retrieved live from PyPI for the failing distribution,
so that every version the model can choose is one that is known to exist.

\section{FAIR Principles and Knowledge Graphs}
\label{sec:rw-fair}

The FAIR Jupyter project this thesis extends is itself grounded in the
FAIR principles: that data, and by extension research software, should
be Findable, Accessible, Interoperable, and
Reusable~\cite{wilkinson2016fair}. Barker et al.~\cite{barker2022fair4rs}
extend these principles specifically to research software. They argue
that the Reusable principle for software requires recording the
\emph{provenance} of its execution environment -- not only the code
itself, but the conditions under which it did or did not run. Samuel
and Mietchen's FAIR Jupyter knowledge
graph~\cite{samuel2024fairkg,samuel2024zenodo} already stores exactly
this kind of provenance for the original execution attempts. Recording
a repair attempt in the same graph, using the PROV
ontology~\cite{provenance2013} for the activity/agent/time structure,
extends that existing provenance record to cover repair as well as
detection. This is the basis for objective O6, designed in
\Cref{chap:architecture}, implemented in \Cref{chap:methodology}, and
validated in \Cref{chap:validation} first on pilot data and then over the
complete table of the final evaluation.

\section{Research Gap}
\label{sec:rw-gap}

Taken together, the literature above shows that the individual building
blocks already exist separately. Pipelines exist that detect and record
reproducibility failures at
scale~\cite{samuel2025docker,samuel2024gigascience,samuel2024fairkg}. Tools
exist that infer or restore a missing Python
environment, from static analysis~\cite{wang2021snifferdog} to
knowledge graphs~\cite{horton2019dockerizeme,ye2022pyego,cheng2022pycre,cheng2024readpye}
to live-retrieval LLMs~\cite{bartlett2025pllm}. LLM agents exist that can fix
errors in notebooks specifically, interactively, one session at a
time~\cite{grotov2024untangling,grotov2024debug}. FAIR-aligned provenance
recording exists for the original failures~\cite{samuel2024fairkg}.

Among the systems reviewed for this thesis, none combines all of the
following: (1) an LLM-based repair
layer that runs autonomously, without a user manually triggering each
attempt; (2) inside an existing reproducibility pipeline rather than on
standalone files; (3) with the repair suggestion grounded in live PyPI data
specifically for the pip-based dependency errors that pipeline already
detects; and (4) paired with a plain-language explanation and, eventually,
provenance recorded as RDF linked to the same knowledge graph the pipeline
already produces.

PLLM~\cite{bartlett2025pllm} is the closest system in mechanism, but it is
not pipeline-embedded, does not explain a failure, and does not write to a
knowledge graph. Grotov et al.'s agents~\cite{grotov2024untangling,
grotov2024debug} are the closest for notebooks specifically, but they assume
an interactive session with a live kernel, not an autonomous batch pass over
a pipeline's own residual failures. \Cref{tab:rw-gap} summarises this
positioning against the thesis objectives introduced in
\Cref{chap:introduction}.

\begin{table}[htbp]
  \centering
  \begin{tabular}{|l|p{4.8cm}|p{4.0cm}|}
    \hline
    \textbf{Objective} & \textbf{Gap it fills} & \textbf{Builds on} \\
    \hline
    O1 -- Explanation & None of the reviewed pipelines explains a residual failure in plain language & Grotov~\cite{grotov2024untangling,grotov2024debug}, Jupyter AI~\cite{jupyterai2023} \\
    \hline
    O2 -- PyPI-grounded repair & None of the reviewed systems embeds PyPI-grounded RAG repair inside a reproducibility pipeline & PLLM~\cite{bartlett2025pllm}, Lewis et al.~\cite{lewis2020rag} \\
    \hline
    O3 -- Fix validation & Autonomous, batch re-execution validation is not done by the reviewed notebook agents & ChatRepair~\cite{xia2024chatrepair}, Grotov~\cite{grotov2024debug} \\
    \hline
    O4 -- Benchmark dataset & No reviewed benchmark covers dependency repair on biomedical notebooks & Grotov~\cite{grotov2024untangling}, Pimentel~\cite{pimentel2019msr} \\
    \hline
    O5 -- Pipeline integration & The FAIR Jupyter pipeline stops at logging; no repair layer exists yet & Samuel et al.~\cite{samuel2025docker,samuel2024gigascience,samuel2024fairkg} \\
    \hline
    O6 -- KG enrichment & None of the reviewed systems records repair provenance as FAIR RDF triples & \cite{samuel2024fairkg}, PROV-O~\cite{provenance2013}, \cite{wilkinson2016fair} \\
    \hline
  \end{tabular}
  \caption{Thesis objectives, the gap each one fills, and the prior work it builds on.}
  \label{tab:rw-gap}
\end{table}

This is the gap the remainder of this thesis addresses: a controlled,
package-registry-grounded LLM repair layer for dependency failures in a FAIR
Jupyter notebook pipeline, that also validates its own repairs by
re-executing the notebook, rather than trusting the model's suggestion on
faith.
